\par\medskip{\large\bfseries\color{mauPhieu} Bài 15. Định lí Thalès trong tam giác}\par\smallskip
\textbf{1. Kiến thức cần nhớ}
\begin{itemize}
\item \textbf{Tỉ số của hai đoạn thẳng} là tỉ số độ dài của chúng (theo cùng một đơn vị đo): $\dfrac{AB}{CD}$. Tỉ số không phụ thuộc đơn vị đo.
\item \textbf{Đoạn thẳng tỉ lệ}: $AB$, $CD$ tỉ lệ với $A'B'$, $C'D'$ nếu $\dfrac{AB}{CD}=\dfrac{A'B'}{C'D'}$.
\item \textbf{Định lí Thalès}: Nếu một đường thẳng song song với một cạnh của tam giác và cắt hai cạnh còn lại thì nó định ra trên hai cạnh đó những đoạn thẳng tương ứng tỉ lệ.\\
Tam giác $ABC$, $B'\in AB$, $C'\in AC$, $B'C'\parallel BC$ $\Rightarrow$ $\boxed{\dfrac{AB'}{AB}=\dfrac{AC'}{AC};\ \dfrac{AB'}{B'B}=\dfrac{AC'}{C'C};\ \dfrac{B'B}{AB}=\dfrac{C'C}{AC}}$.
\item \textbf{Định lí Thalès đảo}: Nếu một đường thẳng cắt hai cạnh của một tam giác và định ra trên hai cạnh này những đoạn thẳng tương ứng tỉ lệ thì đường thẳng đó song song với cạnh còn lại của tam giác.\\
$\dfrac{AB'}{AB}=\dfrac{AC'}{AC}$ (hoặc $\dfrac{AB'}{B'B}=\dfrac{AC'}{C'C}$) $\Rightarrow B'C'\parallel BC$.
\end{itemize}

\textbf{2. Phương pháp giải các dạng}

%% DANG 8H4?15-1
\textbf{Dạng 1. Tỉ số của hai đoạn thẳng, đoạn thẳng tỉ lệ}
\begin{enumerate}
\item Đổi độ dài các đoạn thẳng về cùng một đơn vị đo.
\item Lập tỉ số rồi rút gọn; kiểm tra đoạn thẳng tỉ lệ bằng cách so sánh hai tỉ số (hoặc nhân chéo).
\end{enumerate}
Lưu ý: tỉ số hai đoạn thẳng là số dương, không có đơn vị.

%% DANG 8H4?15-2
\textbf{Dạng 2. Định lí Thalès thuận: tính độ dài}
\begin{enumerate}
\item Chỉ ra đường thẳng song song với một cạnh và cắt hai cạnh còn lại của tam giác.
\item Viết tỉ lệ thức Thalès chứa đoạn cần tìm và các đoạn đã biết.
\item Giải tỉ lệ thức (nhân chéo) để tìm độ dài.
\end{enumerate}
Lưu ý: các đoạn trong một tỉ số phải nằm trên cùng một cạnh và lấy tương ứng (đỉnh chung $A$ ở cả hai vế).

%% DANG 8H4?15-3
\textbf{Dạng 3. Định lí Thalès đảo: chứng minh song song}
\begin{enumerate}
\item Tính hai tỉ số $\dfrac{AB'}{AB}$ và $\dfrac{AC'}{AC}$ (hoặc $\dfrac{AB'}{B'B}$ và $\dfrac{AC'}{C'C}$).
\item Hai tỉ số bằng nhau thì kết luận $B'C'\parallel BC$ theo định lí Thalès đảo.
\end{enumerate}
Lưu ý: hai tỉ số phải lấy tương ứng trên hai cạnh, cùng tính từ đỉnh chung.

%% DANG 8H4?15-4
\textbf{Dạng 4. Ứng dụng thực tế (đo chiều cao, khoảng cách)}
\begin{enumerate}
\item Vẽ hình mô tả, chỉ ra tam giác có hai đoạn thẳng song song (cọc, bóng nắng, cây\dots).
\item Lập tỉ lệ thức Thalès, thay số đo đã biết (cùng đơn vị) để tính độ dài cần tìm.
\item Trả lời bài toán kèm đơn vị.
\end{enumerate}

\par\medskip{\large\bfseries\color{mauPhieu} Bài 16. Đường trung bình của tam giác}\par\smallskip
\textbf{1. Kiến thức cần nhớ}
\begin{itemize}
\item \textbf{Đường trung bình của tam giác} là đoạn thẳng nối trung điểm hai cạnh của tam giác. Mỗi tam giác có ba đường trung bình.
\item \textbf{Định lí}: Đường trung bình của tam giác song song với cạnh thứ ba và bằng nửa cạnh đó.\\
$M$, $N$ lần lượt là trung điểm của $AB$, $AC$ $\Rightarrow$ $\boxed{MN\parallel BC,\ MN=\dfrac12BC}$.
\item \textbf{Nhận xét} (suy từ định lí Thalès): đường thẳng đi qua trung điểm một cạnh của tam giác và song song với cạnh thứ hai thì đi qua trung điểm cạnh thứ ba.
\end{itemize}

\textbf{2. Phương pháp giải các dạng}

%% DANG 8H4?16-1
\textbf{Dạng 1. Tính độ dài nhờ đường trung bình}
\begin{enumerate}
\item Chỉ ra hai trung điểm của hai cạnh $\Rightarrow$ đoạn nối chúng là đường trung bình.
\item Áp dụng $MN=\dfrac12BC$ (hoặc $BC=2MN$) để tính độ dài.
\end{enumerate}
Lưu ý: đường trung bình bằng nửa cạnh \textit{thứ ba} (cạnh không chứa hai trung điểm).

%% DANG 8H4?16-2
\textbf{Dạng 2. Chứng minh trung điểm, song song nhờ đường trung bình}
\begin{enumerate}
\item Chứng minh song song: chỉ ra đoạn thẳng là đường trung bình của một tam giác $\Rightarrow$ song song với cạnh thứ ba.
\item Chứng minh trung điểm: dùng nhận xét — qua trung điểm một cạnh, kẻ song song với cạnh thứ hai thì cắt cạnh thứ ba tại trung điểm.
\item Bài tổng hợp: có thể dùng đường trung bình của nhiều tam giác (ví dụ tứ giác có các trung điểm $\Rightarrow$ hình bình hành).
\end{enumerate}
